\documentclass[11pt]{article}
\usepackage[margin=25mm]{geometry}
\usepackage{fontspec}
\setmainfont{Latin Modern Roman}
\usepackage{amsmath,amssymb,mathtools}
\usepackage{unicode-math}
\setmathfont{Latin Modern Math}
\usepackage{xurl,hyperref}
\hypersetup{colorlinks=true,urlcolor=blue}
\setcounter{secnumdepth}{0}
\setlength{\parindent}{0pt}
\setlength{\parskip}{5pt}
\setlength{\emergencystretch}{3em}
\newcommand{\sref}[2]{\hyperlink{source:#1:#2}{[S#2]}}
\newcommand{\cataloguescope}[1]{\paragraph{Recorded target.}#1}
\newcommand{\R}{\mathbb R}
\newcommand{\C}{\mathbb C}
\newcommand{\N}{\mathbb N}
\newcommand{\Z}{\mathbb Z}
\newcommand{\Q}{\mathbb Q}
\newcommand{\F}{\mathbb F}
\newcommand{\eps}{\varepsilon}
\begin{document}
% UnsolvedMath ID 20002557; source catalogue code AIM-PHYSICS-0065
\setcounter{section}{20002556}
\section{A feasibility sieve for negative-energy binary-single direct scattering}\label{20002557}
{\small UnsolvedMath ID 20002557 · Mathematical Physics}
\subsection{Definitions and mathematical statement}

\paragraph{Shared notation.}$\mathbb N=\{1,2,\ldots\}$; the standard number fields have their usual meanings. All additional conventions are stated locally.


\paragraph{Definitions and notation.}In the Newtonian $n$-body problem, positions $q_i\in\mathbb R^d$ and masses $m_i>0$ obey $m_i\ddot q_i=\nabla_{q_i}U$, where $U(q)=\sum_{i<j}m_im_j/|q_i-q_j|$ and the gravitational constant is normalized to one. Here $T=\tfrac12\sum_i m_i|\dot q_i|^2$, $E=T-U$, and a collision has $q_i=q_j$. The center of mass can be fixed at zero; moment of inertia is $I=\sum_i m_i|q_i|^2$. For a binary plus third body, Jacobi coordinates describe internal binary separation and the third body’s position relative to the binary center of mass. Kepler elements encode an asymptotic two-body conic, including energy, angular momentum, orientation and phase/time conventions. Negative total energy can coexist with an escaping third body when the binary remains bound.

\paragraph{Catalogue wording retained for traceability.}
Consider three-body orbits with negative total energy that are unbound in both time directions. Their Jacobi vectors approach a bound elliptic binary and a hyperbolic escaping orbit in each distant past/future description. Which pairs of incoming and outgoing Kepler elements can be connected by one such orbit?
\paragraph{Scope and source qualification.}Energy and angular-momentum conservation give necessary constraints rather than a full scattering characterization. The source’s asymptotic Kepler description is preserved; phase/time conventions and handling of channel exchange must be specified. The catalogue title/status was paired with a machine-generated research attempt; that attempt is not peer-reviewed evidence of settlement. The statement covers the original source questions rather than only the title’s subcase.
\subsection{Short English statement}
Characterize the incoming and outgoing Kepler data connected by negative-energy three-body scattering.
\subsection{Research attempt}
\paragraph{Model, date and investigated scope.}
OpenAI GPT-6 Astra (configured model \texttt{gpt-6-astra}), ultra reasoning; generated 11 October 2026. Proved the exact spatial orientation interval criterion and an actual equal-mass negative-energy Kepler-data incompatibility example. Scattering sufficiency, phases and channel exchange remain uncharacterized.

\paragraph{Sources and literature check.}
Read the supplied complete definition, catalogue wording, scope qualifications and source list for this canonical ID. The original-source citations are inherited as supplied references; no fresh primary-source or novelty search was performed for this attempt. No current open-status claim is inferred from this calculation.

\paragraph{An orientation-feasibility obstruction.}For a fixed incoming channel, let the magnitudes of the binary and escaping Kepler angular momenta be \(L_b^-,L_o^-\ge0\); define outgoing magnitudes similarly. If spatial orientations are unspecified, the possible magnitude of their vector sum is exactly
\[I^\pm=[|L_b^\pm-L_o^\pm|,L_b^\pm+L_o^\pm].
\]
The triangle inequality gives necessity. Sufficiency follows by choosing the angle \(\theta\) between vectors: the squared total is \(L_b^2+L_o^2+2L_bL_o\cos\theta\), which fills the entire interval as \(\theta\) varies. If one magnitude is zero, the interval is a singleton and the same conclusion holds. Conservation therefore requires \(I^-\cap I^+\ne\varnothing\), even after allowing arbitrary orientations.
\paragraph{An equal-mass three-body obstruction.}Take all three masses equal to one and use the same binary channel on both ends. The binary Kepler Hamiltonian has reduced mass \(\mu_b=1/2\) and gravitational product \(k_b=1\); the outer binary-single Hamiltonian has \(\mu_o=2/3\), \(k_o=2\). For \(p^2/(2\mu)-k/r\), energy e and angular momentum L give eccentricity squared \(1+2eL^2/(\mu k^2)\), by the radial turning-point quadratic. Choose binary energy \(-1\), binary angular magnitude \(1/4\), and outer energy \(1/2\) on both ends. The total energy is \(-1/2\), and the binary eccentricity squared is \(3/4\). Set incoming outer angular magnitude \(1/4\) and outgoing magnitude \(2\). Their outer eccentricities squared are respectively \(1+3/128\) and \(5/2\), so both are valid nondegenerate hyperbolic Kepler data. Nevertheless
\[I^-=[0,1/2],\qquad I^+=[7/4,9/4],\]
which are disjoint. These energy-matched asymptotic elements cannot be joined by any three-body orbit, regardless of orientation or phase choices.
\paragraph{Gap.}Energy/angular-momentum compatibility is only necessary. Phases, orientations already fixed in the data, collision avoidance, time conventions and channel exchange impose further restrictions. Constructing the actual scattering relation for specified three-body masses remains the first missing dynamical step. The explicit example uses the actual reduced masses of one equal-mass three-body channel.

\paragraph{Verification and adversarial audit.}
Checked triangle sufficiency, zero-magnitude endpoints, actual equal-mass Jacobi reduced masses and gravitational products, both hyperbolic eccentricities, elliptic admissibility and disjoint conserved angular-momentum intervals. This was an LLM adversarial audit, not independent human verification; no theorem prover was used. Prompts A then B applied: replaced an illustrative normalized-channel example by actual equal-mass three-body parameters and rechecked every eccentricity and angular interval. The condition remains necessary only; constructing the full scattering relation is unproved.

Coordinator checked the attainable angular-momentum interval and all reduced masses, energies and eccentricities in the equal-mass obstruction example. This additional check was model-assisted, not independent human verification.

\paragraph{Tools, process and contributions.}
Python 3 and local shell used to read supplied files and serialize this manuscript; mathematics derived symbolically in the text, with no numerical experiment or formal proof assistant. Contribution prompt A was followed by prompt B. GPT-6 Astra with ultra reasoning drafted the research and performed its stated self-audit; any additional coordinator contribution is identified in the verification and audit. A GPT-6-based Codex coordinator assembled the manuscript and checked repository integration; its exact variant and reasoning configuration were not exposed. No human mathematical contribution, independent human proof check or formal proof certification is supplied. Complete derivations and the consulted source references are included here; private reasoning transcripts are not published.

\paragraph{Final status and first remaining gap.}
\textbf{UNRESOLVED}. Construct the scattering map/relation for actual channel reduced masses and prove which compatible asymptotic data are dynamically connected.

\section{Completion Estimate}
\noindent\textbf{COMPLETION: 0\%}
The calculation is an elementary benchmark or conditional reduction; it establishes no claimed new portion of the complete source target.

\subsection{Sources}
\begin{itemize}
\item[\hypertarget{source:20002557:1}{S1}] ULAM.AI and the UnsolvedMath Contributors, UnsolvedMath, record 20002557 (AIM-PHYSICS-0065). Catalogue curation under CC BY 4.0; definitions and English exposition added locally.
\url{https://huggingface.co/datasets/ulamai/UnsolvedMath}.
\item[\hypertarget{source:20002557:2}{S2}] American Institute of Mathematics, Variational Methods in Celestial Mechanics; notes by Kuo-Chang Chen, version 22 June 2003, source item 28.
\url{https://aimath.org/WWN/varcelest/varcelest.pdf}.
\end{itemize}
\label{end:20002557}
\end{document}
